% Generated table; it uses the manuscript preamble's macros (\hltint, \hcell, \altrow, PanelBg).
\begin{xltabular}{\textwidth}{>{\raggedright\arraybackslash}X r r r r r r}
\caption{\hltint{HlM}{Alternative spatial splits and the severity mix by area.}}\label{tab:a_spatialalt}\\
\toprule
\headerrow
\hcell{Split} & \hcell{$n$ first} & \hcell{$n$ second} & \hcell{Covariates} & \hcell{LR} & \hcell{d.f.} & \hcell{$p$ boot.} \\
\midrule
\endfirsthead
\multicolumn{7}{l}{\cellcolor{HeaderGold}\textcolor{HeaderText}{\textbf{\tablename\ \thetable{} \textit{(continued)}}}} \\
\toprule
\endhead
\midrule
\multicolumn{7}{r}{\textit{Continued on next page}} \\
\endfoot
\bottomrule
\multicolumn{7}{p{0.97\linewidth}}{\footnotesize \textit{Note:} With two and six of 35 covariates estimable on both sides, these tests do not test the M1 specification, and they are reported as descriptive checks. In Panel~A the first group holds the records inside the urban area, or the rural records. The six major cities take precedence in Panel~B, and the rural group holds the remaining records with the rural flag set.}\\
\endlastfoot
\addlinespace
\rowcolor{PanelBg}
\multicolumn{7}{l}{\textbf{Panel A. Alternative spatial splits}} \\
\addlinespace[2pt]
Inside against outside an urban area of 200,000 or more & 404 & 118 & 6 & 7.30 & 8 & 0.511 \\
\altrow
Rural against non-rural & 63 & 459 & 2 & 0.75 & 4 & 0.938 \\
\addlinespace
\rowcolor{PanelBg}
\multicolumn{7}{l}{\textbf{Panel B. Severity mix of three spatial groups}} \\
\addlinespace[2pt]
\midrule
\headerrow
\hcell{Group} & \hcell{$n$} & \hcell{O (\%)} & \hcell{BC (\%)} & \hcell{KA (\%)} & \hcell{} & \hcell{} \\
\midrule
Six major cities & 257 & 16.7 & 70.8 & 12.5 & & \\
\altrow
Other urban & 202 & 12.9 & 67.8 & 19.3 & & \\
Rural & 63 & 12.7 & 69.8 & 17.5 & & \\
\end{xltabular}
